\documentclass[11pt]{article}

\usepackage[margin=1in]{geometry}
\usepackage{microtype}
\usepackage{amsmath}
\usepackage{graphicx}
\usepackage{booktabs}
\usepackage[hidelinks]{hyperref}

\graphicspath{{figures/}}
\setlength{\emergencystretch}{1.2em}

\newcommand{\natex}{\texttt{natex}}

\title{Still No Kink at o1: A Precision-Weighted Rerun of the Epoch
Capabilities Index on a Fresh Vintage}
\author{Hauke Hillebrandt\thanks{University College London. Email:
\texttt{ucjthhi@ucl.ac.uk}. This paper and the underlying \natex{} software were
prepared with substantial assistance from Anthropic's Claude models; the author
reviewed the analyses and text and is responsible for all remaining errors.}}
\date{July 2026 \\ \small analysis pass of record: 2026-07, \natex{} v0.2.0, seed 0}

\begin{document}

\maketitle

\begin{abstract}
\noindent The aggregate AI capabilities trend did not bend at o1. On a fresh
vintage of the Epoch Capabilities Index (ECI; 211 models, February 2023 --
July 2026, twelve more months than the prior analysis pass), a sharp
regression kink design in calendar time at the o1-preview release (12
September 2024) estimates an all-models slope change of $+0.0065$ ECI
points/day (HC1 se $0.0131$, $t=0.50$, $p=0.618$, 95\% CI
$[-0.019, +0.032]$; slopes $\approx22.3$ pre vs.\ $\approx24.7$ points/yr
post), null in 17 of 18 specification cells --- while the same specification
rejects at 4 of 8 placebo cutoffs away from the release, so the design has
power and the dial reads zero. A per-model precision-weighted rerun (weights
$1/\sigma_i^2$ from the published confidence intervals) returns a nominally
significant \emph{negative} kink ($-0.0190$, se $0.0079$, $t=-2.41$,
$p=0.016$; 18/18 cells), but triage identifies it as an artifact of smooth
post-2024 concavity rather than an o1-localized break: placebo cutoffs at
$+90$/$+180$/$+270$ days give same-sign, larger-or-comparable kinks
($-0.0399$, $p=0.0064$; $-0.0297$, $p=5.5\times10^{-6}$; $-0.0154$,
$p=0.021$), and a local-quadratic fit absorbs it ($p=0.140$ at bandwidth
540). The clean frontier-records series (24 records, only 4 pre-cutoff
points in the estimation cell) remains a non-estimable design: its nominal
$t=2.28$ dies under a 30-day donut ($t=0.07$), flips sign under a 45-day
donut and at bandwidth 730, and its placebo grid rejects 5/8. A
high-confidence subset (excluding 26 Unverified/Speculative models)
reproduces every conclusion, and there is no release-density kink at the
cutoff ($t=-0.42$). The prior-vintage null replicates and is upgraded; per
the collection's ranking specification this result folds into the flagship
paper's ECI section as an appendix rather than standing as an independent
finding.
\end{abstract}

\section{Introduction}\label{sec:intro}

OpenAI released o1-preview on 12 September 2024, the first widely deployed
``reasoning model'' trained with large-scale reinforcement learning to spend
inference-time compute on a chain of thought \cite{openai2024}. The release
is a natural candidate for a regime change in AI capabilities: contemporaneous
research argued that scaling test-time compute can beat scaling parameters
\cite{snell2024}, and the reasoning-model recipe was adopted across frontier
labs within months. If the recipe changed the \emph{rate} of capability
progress, an aggregate capability index should show a slope change --- a kink
--- in calendar time at or shortly after the release.

Measuring that requires a scale on which models years apart are comparable.
The Epoch Capabilities Index (ECI) stitches dozens of benchmarks onto a single
latent-ability scale via an item-response-theory model, precisely to support
longitudinal comparisons of this kind \cite{ho2025}. This paper asks a narrow
question: does the ECI trend kink at o1?

The estimand --- a change in trend slope at a known date --- is the regression
kink design (RKD) of \cite{card2015}, applied here with calendar time as the
running variable; \cite{bockerman2025} formalize the closely related
difference-in-kinks design that the same toolkit implements. Two strands of
methodological literature discipline the exercise. First, regression
discontinuity (and kink) designs in \emph{time} lack the randomization logic
of cross-sectional designs: units near the cutoff are not quasi-randomly
assigned, and serial dependence and smooth trends generate spurious breaks
\cite{hausman2018}. Second, kink estimates on smooth series are known to
produce spuriously significant slope changes, which is why placebo
distributions at non-event dates are the appropriate yardstick
\cite{ganong2018}. Both concerns are operationalized below through a
mandatory placebo-cutoff battery and curvature checks, estimated with the
\natex{} toolkit \cite{natex}.

This note is a robustness upgrade of a prior null: an earlier analysis pass on
a 2025 vintage of the index found no kink at o1. The fresh vintage adds
twelve months of post-cutoff data and per-model confidence intervals, which
enable a precision-weighted re-estimate --- the natural next test, since noisy
early-2023 models could in principle mask a true kink in the unweighted fit.

\section{Data}\label{sec:data}

The outcome is the fresh vintage of Epoch AI's \texttt{eci\_scores.csv}
(fetched July 2026): 211 aggregated models with release dates from 2023-02-24
to 2026-07-09 and no missing score or date, spanning ECI $\approx55$ to
$161.77$. The vintage adds roughly twelve months of data relative to the
prior analysis pass. Each model carries a 95\% confidence interval; the
per-model precision weight is $w_i = 1/\sigma_i^2$ with
$\sigma_i = (\mathrm{ci}_{\mathrm{high},i} - \mathrm{ci}_{\mathrm{low},i})/3.92$.
Two models (Claude 3.5 Sonnet and GPT-5, the scale's two anchor models) have
no published CI; their $\sigma_i$ is imputed at the sample median ($2.005$)
for the weighted runs only.

A version-level Confidence flag is attached by a two-pass name join against
Epoch's per-version index file; 189 of 211 models match (the 22 unmatched are
treated as not-flagged), and 26 models (12.3\%) are flagged
Unverified/Speculative --- including the current \#1, GPT-5.6 Sol (ECI
$161.77$), and GPT-5. A high-confidence subset excludes these 26 models
($n=185$).

Two series are analyzed. The \emph{all-models} series uses all 211 releases.
The \emph{clean frontier-records} series keeps strict running-maximum records
of the score (computed over non-missing scores only, in date order, with a
stable same-day tie-break): 24 records, 7 pre-cutoff and 17 post-cutoff ---
and only 4 pre-cutoff points inside bandwidth 540 days, 3 after a 45-day
donut, which already flags it as a marginal design.

\section{Design and Methods}\label{sec:methods}

\paragraph{Sharp RKD in calendar time.} The running variable is days since
the o1-preview release (cutoff $= 0$ at 2024-09-12). The estimand is the
change in the slope of ECI in points/day at the cutoff, estimated by local
linear regression on each side within the bandwidth, triangular kernel, HC1
standard errors (\natex{} v0.2.0 \cite{natex}, \texttt{natex kink} with a
unit policy-kink change; \cite{card2015,bockerman2025}). The primary
specification uses bandwidth 540 days, donut 0 ($n=179$; 85 left, 94 right).
A specification grid varies bandwidth $\in\{365,540,730\}$ days $\times$
kernel $\in\{\text{triangular},\text{uniform}\}$ $\times$ donut
$\in\{0,30,45\}$ days
(18 cells per series--weighting pair). The CI-weighted estimator is a local
WLS that mirrors \natex{} exactly (same selection, kernel, block design, HC1)
with weights kernel$\,\times\,1/\sigma_i^2$; with unit weights it reproduces
\texttt{natex regression\_kink} in all 68 evaluable grid cells (tolerance
$10^{-8}$ on points, $10^{-6}$ on SEs) and matches the NaN reasons in the
rest, and the CLI headline cells match the grid to all printed digits. The
design is deterministic; seed 0 is declared and no stochastic step exists.

\paragraph{Honest-inference caveats, stated as run.} This is a calendar-time
RKD: model releases are not randomized around the cutoff, so the
cross-sectional RD identification logic does not apply and smooth trends can
masquerade as kinks \cite{hausman2018}. HC1 treats observations as
independent, ignoring same-day multi-releases and lab-level clustering. The
placebo-cutoff grid --- the same specification re-estimated at $\pm90$,
$\pm180$, $\pm270$, $\pm360$ days, in the spirit of \cite{ganong2018} --- is
therefore the operative guard, and below it is what both establishes power
for the null and disqualifies the weighted ``significant'' cell. A
local-quadratic (degree-2) refit checks whether any kink survives allowing
smooth curvature; a release-density kink test (binned 30-day release counts)
checks for composition churn at the cutoff.

\section{Results}\label{sec:results}

\begin{figure}[t]
\centering
\includegraphics[width=\textwidth]{fig1}
\caption{(a) The 211 fresh-vintage ECI scores against days since the
o1-preview release (vertical line), with per-side local-linear fits
(bandwidth 540 days, triangular kernel): unweighted (blue) and CI-weighted
(orange). The unweighted slopes are statistically indistinguishable across
the cutoff; the CI-weighted fit bends down. (b) Kink estimates with 95\% HC1
confidence intervals at the o1 cutoff and at eight placebo cutoffs, both
weightings (points offset horizontally for legibility); filled markers are
nominally significant at the 5\% level. The unweighted design rejects at 4/8
placebo cutoffs while reading zero at o1; the CI-weighted negative kink at
o1 is matched by same-sign, larger placebos at $+90$ and $+180$ days.}
\label{fig:fig1}
\end{figure}

\paragraph{The unweighted null, now with more power.} The all-models headline
kink (Table~\ref{tab:kink}, first row) is $+0.00653$ ECI points/day (se
$0.01308$, $t=0.50$, $p=0.618$, 95\% CI $[-0.0191, +0.0322]$): the trend runs
$0.0611$ points/day before o1 and $0.0677$ after ($\approx22.3$ vs.\
$\approx24.7$ points/yr). Of the 18 grid cells, 17 are null; the single
nominal rejection (bandwidth 365, triangular, donut 30: $-0.0547$, se
$0.0269$, $t=-2.03$, $p=0.042$) is the grid's smallest cell and negative.
The placebo grid shows the null is not for lack of power: at bandwidth 540
the same specification rejects at 4 of 8 non-event cutoffs ($-180$ days:
$+0.0401$, $p=0.023$; $+180$: $-0.0323$, $p=7.1\times10^{-4}$; $+270$:
$-0.0288$, $p=1.0\times10^{-4}$; $+360$: $-0.0247$, $p=0.0034$) while the
true cutoff is null in all cells --- the instrument moves elsewhere and reads
zero at o1. The prior-vintage null replicates on twelve more months of data.

\paragraph{The CI-weighted ``kink'' is a curvature artifact.} Precision
weighting flips the headline to a nominally significant \emph{negative} kink:
$-0.01903$ (se $0.00788$, $t=-2.41$, $p=0.0158$), significant in 18 of 18
grid cells ($|t|$ from $2.12$ to $3.32$), with slopes $0.0522$ pre vs.\
$0.0331$ post ($\approx19.0$ vs.\ $\approx12.1$ points/yr). Two facts
disqualify an o1-localized reading. First, the placebo grid rejects with the
\emph{same sign and larger or comparable magnitude} away from the release:
$-0.0399$ ($p=0.0064$) at $+90$ days, $-0.0297$ ($p=5.5\times10^{-6}$) at
$+180$, $-0.0154$ ($p=0.021$) at $+270$ (empirical size 3/8 at
$\alpha=0.05$). Second, a local-quadratic fit absorbs it: degree-2 kinks are
$-0.0392$ ($t=-1.48$, $p=0.140$) at bandwidth 540 and $-0.0381$ ($t=-1.74$,
$p=0.081$) at bandwidth 730. The precision-weighted index decelerates
smoothly through 2025 --- a concavity, not a kink at o1.

\paragraph{The frontier-records series is still not a design.} With only 4
left-cell points at bandwidths $\leq540$ (3 after donut 45), the records
series produces a nominal positive kink at the headline specification
($+0.01313$, se $0.00576$, $t=2.28$, $p=0.023$) that no perturbation
survives: a 30-day donut kills it ($t=0.07$, $p=0.943$), a 45-day donut flips
it to significantly \emph{negative} ($-0.02817$, $t=-2.52$, $p=0.012$),
bandwidth 730 flips the sign ($-0.00429$, $t=-0.27$), and the placebo grid
rejects 5/8 --- with the $-360$/$-270$ cells exploding to $-1.14$
($p=3\times10^{-4}$) on $n=11$. The CI-weighted records headline is null
($+0.00794$, $t=1.31$, $p=0.191$). The fresh vintage does not clear the
left-cell bar: the series remains non-estimable, exactly as in the prior
pass.

\paragraph{Robustness and diagnostics.} Excluding the 26
Unverified\slash Speculative models ($n=185$) changes nothing: unweighted
$+0.01248$ ($t=0.88$, $p=0.379$, null), CI-weighted $-0.02109$ ($t=-2.32$,
$p=0.020$) with the same placebo failure (weighted empirical size 4/8,
unweighted 6/8). The release-density check finds no composition-churn kink
at the cutoff: binned 30-day release counts give $-0.00410$ (se $0.00967$,
$t=-0.42$, $p=0.672$), an improvement on the prior vintage's marginal
$t=1.82$.

\begin{table}[t]
\centering
\small
\caption{Sharp RKD in time at o1-preview (2024-09-12), ECI points/day.
Headline specification: bandwidth 540 days, triangular kernel, donut 0, HC1.}
\label{tab:kink}
\setlength{\tabcolsep}{4.5pt}
\begin{tabular}{llrrrrr}
\toprule
Series & Weighting & Kink & se & $t$ & $p$ & $n$ \\
\midrule
\multicolumn{7}{l}{\emph{Panel A: headline cells}}\\
All models (211)            & unweighted  & $+0.00653$ & $0.01308$ & $0.50$  & $0.618$ & 179 \\
All models                  & CI-weighted & $-0.01903$ & $0.00788$ & $-2.41$ & $0.016$ & 179 \\
Frontier records (24)       & unweighted  & $+0.01313$ & $0.00576$ & $2.28$  & $0.023$ & 18 \\
Frontier records            & CI-weighted & $+0.00794$ & $0.00608$ & $1.31$  & $0.191$ & 18 \\
High-confidence (185)       & unweighted  & $+0.01248$ & $0.01420$ & $0.88$  & $0.379$ & 162 \\
High-confidence             & CI-weighted & $-0.02109$ & $0.00908$ & $-2.32$ & $0.020$ & 162 \\
Release density (30-d bins) & ---         & $-0.00410$ & $0.00967$ & $-0.42$ & $0.672$ & 36 \\
\midrule
\multicolumn{7}{l}{\emph{Panel B: triage of the nominally significant cells}}\\
All models, CI-w., placebo $+90$ d  & CI-weighted & $-0.03986$ & $0.01462$ & $-2.73$ & $0.0064$ & 181 \\
All models, CI-w., placebo $+180$ d & CI-weighted & $-0.02966$ & $0.00653$ & $-4.54$ & $5.5\times10^{-6}$ & 174 \\
All models, CI-w., degree 2, bw 540 & CI-weighted & $-0.03922$ & $0.02659$ & $-1.48$ & $0.140$ & 179 \\
Frontier records, donut 30          & unweighted  & $+0.00048$ & $0.00667$ & $0.07$  & $0.943$ & 16 \\
Frontier records, donut 45          & unweighted  & $-0.02817$ & $0.01118$ & $-2.52$ & $0.012$ & 15 \\
Frontier records, bw 730            & unweighted  & $-0.00429$ & $0.01592$ & $-0.27$ & $0.788$ & 24 \\
\bottomrule
\end{tabular}

\medskip
\parbox{0.95\textwidth}{\footnotesize Notes: kinks in ECI points/day; HC1
standard errors against standard normal critical values. CI weights
$1/\sigma_i^2$, $\sigma_i = (\mathrm{ci}_{\mathrm{high}} -
\mathrm{ci}_{\mathrm{low}})/3.92$. Placebo-cutoff grids (8 cutoffs at
$\pm90/180/270/360$ days, bandwidth 540, triangular): all-models unweighted
rejects 4/8 (true cutoff null in all cells); all-models CI-weighted rejects
3/8, same sign as the headline; frontier records rejects 5/8. Panel B donut
and bandwidth rows are the perturbations of the frontier-records headline
cell.}
\end{table}

\section{Caveats and Conclusion}\label{sec:conclusion}

The caveats of record are these. First, this is a calendar-time RKD: model
releases are not randomized around the cutoff, and HC1 treats them as
independent (same-day multi-releases and lab-level clustering are ignored);
the placebo-cutoff grid is the operative guard, and it is what both
establishes power for the null and disqualifies the weighted ``significant''
cell. Second, data quality: 26 of 211 models (12.3\%) carry
Unverified/Speculative version-level Confidence flags --- including the
current \#1, GPT-5.6 Sol (161.77), and GPT-5 --- and the confidence join
matched only 189/211 models (22 unmatched treated as not-flagged); the
high-confidence robustness subset reproduces every conclusion. Third, two
models (Claude 3.5 Sonnet and GPT-5) have no published CI and enter the
weighted runs at the median imputed $\sigma$; and the index itself is a
modeled latent scale whose vintages revise as benchmarks are added
\cite{ho2025}, which is why the replication across vintages matters.

The verdict is an upgraded null with power. On twelve more months of data,
the aggregate ECI trend shows no kink at o1: the unweighted estimate is a
precise zero at a cutoff where the design demonstrably has power to reject
(4/8 placebos reject elsewhere), the one new nominally significant result ---
the CI-weighted negative kink, significant in 18/18 cells --- is correctly
identified by triage as a smooth post-2024 deceleration of the
precision-weighted index (same-sign larger placebos at $+90$/$+180$,
absorbed by a local quadratic) rather than an o1-localized break, and the
genuine frontier-records series remains a four-left-point non-design.
Whatever the reasoning-model recipe did to AI capabilities, it did not bend
the aggregate index's calendar-time trend at the release date. Per the
collection's ranking specification, this result folds into the flagship
paper's ECI section as an appendix rather than standing as an independent
mini-paper.

\paragraph{Reproducibility.} All estimates were produced with \natex{} v0.2.0
\cite{natex}; estimation is deterministic (local weighted least squares),
seed 0 was declared, and no stochastic step exists in the design. The
underlying \texttt{eci\_scores.csv} vintage is not committed to the
repository; the run of record retains the derived panels, the 144-cell grid,
the placebo grids, and the validation counts. Figure~\ref{fig:fig1}
regenerates deterministically from the committed
\texttt{figures/make\_fig.py}, which asserts the headline estimates, the
recomputed fit slopes, and the placebo rejection counts against the numbers
of record before drawing.

\footnotesize
\begin{thebibliography}{9}
\setlength{\itemsep}{2pt}\setlength{\parskip}{0pt}

\bibitem{bockerman2025}
B\"ockerman, P., Jysm\"a, S., and Kanninen, O. (2025).
\emph{Difference-in-Kinks Design}.
IZA Discussion Paper No.~18313.
\url{https://docs.iza.org/dp18313.pdf}

\bibitem{card2015}
Card, D., Lee, D.~S., Pei, Z., and Weber, A. (2015).
Inference on causal effects in a generalized regression kink design.
\emph{Econometrica}, 83(6), 2453--2483.

\bibitem{ganong2018}
Ganong, P., and J\"ager, S. (2018).
A permutation test for the regression kink design.
\emph{Journal of the American Statistical Association}, 113(522), 494--504.

\bibitem{hausman2018}
Hausman, C., and Rapson, D.~S. (2018).
Regression discontinuity in time: Considerations for empirical applications.
\emph{Annual Review of Resource Economics}, 10, 533--552.

\bibitem{ho2025}
Ho, A., Denain, J.-S., Atanasov, D., Albanie, S., and Shah, R. (2025).
A Rosetta Stone for AI benchmarks.
arXiv:2512.00193.

\bibitem{natex}
Hillebrandt, H. (2026).
\emph{natex: automated natural-experiment discovery and estimation} (version
0.2.0). Software.
\url{https://github.com/HaukeHillebrandt/natex}

\bibitem{openai2024}
OpenAI (2024).
OpenAI o1 System Card.
arXiv:2412.16720.

\bibitem{snell2024}
Snell, C., Lee, J., Xu, K., and Kumar, A. (2024).
Scaling LLM test-time compute optimally can be more effective than scaling
model parameters.
arXiv:2408.03314.

\end{thebibliography}

\end{document}
